\section{Discussion}
\label{sec:discussion}

\subsection{The Cliff Is About Residual Poison}
The grid and the partial-deletion block agree on a simple statement: what
matters is whether any poison remains, not how much was removed. A model trained
on 2{,}250 of 2{,}500 poison samples behaves like a fully poisoned model; a model
trained on 125 behaves like a nearly clean one. The transition is sharp enough
that the difference between a detector with 90\% recall and one with 100\% recall
is the difference between no defence and a complete defence. This is why the
success tolerance of deletion is exactly $\rstar{=}1.00$ rather than a value
slightly below one.

The result is an empirical statement about this family of attacks and models, not
a theorem. A trigger with a smaller effective support, a lower poison rate, or a
model with more capacity might move the step. But the step was in the same place
across two architectures, two triggers, and two modalities, and it was visible
from both directions (increasing recall in the grid, increasing residual poison
in E2). We would be surprised if it moved far under moderate changes to the
protocol.

\subsection{Detector Engineering Should Optimise Recall}
The precision axis is nearly flat. At fixed recall, adding clean false positives
to the forget set changes ASR by less than a point for deletion and by a small,
non-monotone amount for the recipe. The cost appears on the utility side, where
the recipe loses at most about 0.6 points across the precision cells. This
suggests a division of labour: detector design should maximise recall subject to
a precision floor that keeps the forget set affordable, and removal design should
assume that the forget set will contain clean data.

The activation-clustering result at 10\% poison is a warning about detector
design as much as about removal. The detector returns nothing at exactly the
poison rate where the attack is strongest, because its cluster-size threshold
assumes that poison is a small minority. A detector that cannot flag a large
poison fraction is not conservative; it is silently permissive. Reporting recall
and precision at a single poison rate would not have revealed this.

\subsection{Why Trigger-Aware Repair Tolerates Half the Poison}
The recipe does not need to identify every poisoned sample because its
inoculation stage changes the trigger-to-target association for the whole
network. Label repair removes the label noise that would otherwise re-teach the
association on the flagged subset; inoculation teaches the network that the
trigger is uninformative on a large sample of clean data. The two mechanisms
cover different failures: repair handles the poison that was found, inoculation
handles the poison that was missed, up to a point. The balance breaks at
$r{=}0.25$, where three-quarters of the poison keeps its original label and the
inoculation signal is not enough.

This also explains why the recipe's clean-accuracy cost grows with architecture
depth and trigger strength. Inoculation is a data-dependent regulariser, and the
network must fit both the retained data and the re-labelled trigger examples. The
cost is small (about one point) on ResNet-18 and larger (up to 4.1 points) on
VGG-16-BN with Blend, which is the setting with the largest distribution shift
between clean and trigger-stamped inputs.

\subsection{The Known-Trigger Bound Is the Central Limitation}
The recipe's tolerance is bought with the true trigger, and the synthesis
experiment shows that the purchase cannot be avoided by estimating the trigger
from the model. On both CIFAR attacks the synthesised mask had zero or near-zero
overlap with the true patch, and inoculation with it left ASR above 95\%. This
places the recipe on the ceiling side of the problem: it is a strong upper bound
on what post-training repair can do, and it is not a blind defence.

The bound matters for how the paper should be cited. A reader should not take
away ``repair removes backdoors at 50\% recall''. The correct takeaway is
``repair removes backdoors at 50\% recall if the trigger is known; if it is not,
no method in this study does''. The first statement is a benchmark result. The
second is the deployment result.

\subsection{Relation to Prior Work}
Our findings are consistent with the negative result of Pawelczyk \etal{} for
approximate unlearning \cite{pawelczyk}, and they add a controlled curve that
separates the effect of missed poison from the effect of deleted clean data. They
are also consistent with the relearning result of Hu \etal{} \cite{jogging}: the
oracle-deleted model looks clean by membership inference and by output
divergence, yet relearns the backdoor in ten steps. The difference is that our
probe is a vision classifier and the poison is a trigger, not a text corpus, so
the dormancy is not a language-model artefact.

The ascent artefact is, as far as we know, under-reported. Evaluations of
unlearning that use feature-space ascent without a normalisation control may be
measuring batch-norm corruption. We recommend that any ascent-based removal
result be reported alongside a run with batch-norm statistics frozen to the
retain set.

\subsection{A Procedure for Detector-Bounded Removal}
The results support a short procedure for anyone who has to remove a backdoor
after training, and each step is tied to a measurement in this paper.

\emph{1. Measure recall on ground-truth poison, not detection F1.} The quantity
that predicts removal outcome is the fraction of poison in the forget set. A
detector report that quotes only F1 or only a threshold hides it. A small
poisoned pilot with known poison labels gives the recall curve directly and costs
a fraction of the removal runs it will govern.

\emph{2. Without trigger knowledge, treat deletion as all-or-nothing.} In our
grid the smallest tested recall at which retraining on the retained set passes
the joint criterion is $1.00$; at $0.95$ the ASR is low but the clean-accuracy
cost already fails the criterion, and at $0.90$ the backdoor is at full strength.
A detector whose upper confidence bound on recall does not reach the top of that
cliff should not be used to justify a deletion-based removal claim.

\emph{3. With trigger knowledge, repair is an option, with a budget.} The recipe
holds ASR below $5\%$ down to $r{=}0.5$ at precision 1.0 and below $3\%$ at
precision 0.7, at a clean-accuracy cost of about $1.3$ points on the ResNet-18
grid; on the transferred settings of Section~\ref{sec:general} the cost rises to
$4.09$ points in the hardest case (VGG-16-BN with Blend at $r{=}0.5$). The
budget, not the average, is what a deployment decision needs.

\emph{4. Validate with relearning, and control for batch norm.} A removal claim
that stops at ASR and membership inference can be satisfied by suppression: our
oracle-deleted model returns to $97.7\%$ ASR within ten steps. Any ascent-based
result should be reported next to a run with batch-norm statistics frozen, since
without that control the apparent removal may be normalisation damage.

Nothing in this procedure requires a new method. It requires reporting the
detector's recall, the precision at which it was measured, whether the trigger
was known, and a relearning probe. Those four items are enough to make removal
results comparable across papers.

\subsection{What Would Change the Conclusion}
Three developments would move the conclusion from negative to positive. A
detector with recall at or near 1.0 at usable precision would make deletion
viable, and the cliff would become irrelevant. A trigger-synthesis method with
high IoU on realistic attacks would make the recipe blind, and its tolerance
would become deployable. An amplification signal that survives matched controls
and a utility budget would make iterative detect-remove loops viable. None of
these was observed here. Until one is, the honest engineering advice is to invest
in provenance and duplicate detection before training, because post-training
removal cannot compensate for a detector that misses poison.
